% !PRINTED 28
\weberpage{28}

Das Produkt des konstanten Faktors auf der linken Seite ist die
Diskriminante $D$, und man erhält also nach (13)
\weq{21}{\begin{aligned}
D\,y_1^2y_2^2y_3^2y_4^2 &= f(\Phi, -\Psi)\\
&= a_0\Phi^4 - a_1\Phi^3\Psi + a_2\Phi^2\Psi^2 - a_3\Phi\Psi^3 + a_4\Psi^4,
\end{aligned}}
und die $y_i^2$ sind also die linearen Faktoren dieser biquadratischen
Form.

\begin{center}
\rule{2cm}{0.4pt}
\end{center}

Es gibt noch eine weitere simultane Kovariante der Form $\varphi$,
$\psi$, nämlich die \gesperrt{Jacobische} Funktionaldeterminante der vier
Funktionen $\varphi$, $\psi$, $\Phi$, $\Psi$. Wir bezeichnen sie so:
\weq{22}{K = \begin{vmatrix}
\varphi_1, & \varphi_2, & \varphi_3, & \varphi_4\\
\psi_1, & \psi_2, & \psi_3, & \psi_4\\
\Phi_1, & \Phi_2, & \Phi_3, & \Phi_4\\
\Psi_1, & \Psi_2, & \Psi_3, & \Psi_4
\end{vmatrix}.}
Bildet man sie für die kanonische Form, so ergibt sich nach (12),
(14), (15)
\[
K = -\Delta\,y_1y_2y_3y_4,
\]
und da $\Delta^2 = D$ ist:
\[
K^2 = D\,y_1^2y_2^2y_3^2y_4^2,
\]
also haben wir nach (21)
\weq{23}{K_2 = f(\Phi, -\Psi).} % [sic] printed K_2; derivation requires K^2

\begin{center}
\rule{2cm}{0.4pt}
\end{center}

Die biquadratische Form $f(\Phi, -\Psi)$ hat zwei Kovarianten, die
wir wie in Bd. I, §\,70 so bezeichnen:
\weq{24}{\begin{aligned}
H &= \frac{1}{3}\begin{vmatrix}
\dfrac{\partial^2 f}{\partial \Phi^2}, & -\dfrac{\partial^2 f}{\partial \Phi\,\partial \Psi}\\
\dfrac{\partial^2 f}{\partial \Psi\,\partial \Phi}, & -\dfrac{\partial^2 f}{\partial \Psi^2}
\end{vmatrix},\\
T &= \frac{1}{12}\begin{vmatrix}
f'(\Phi), & f'(-\Psi)\\
H'(\Phi), & H'(-\Psi)
\end{vmatrix}.
\end{aligned}}
Sie sind vom vierten und sechsten Grade in der $\Phi$, $\Psi$, also
vom achten und zwölften Grade in den $x_i$. Es sind \gesperrt{Kovarianten der
Grundkurve} (Kombinanten von $\varphi$ und $\psi$), und es besteht
zwischen ihnen die Relation:
\weq{25}{H^3 - 48\,A\,H\,f^2 - 64\,B\,f^3 = -27\,T^2.}

\begin{center}
\rule{2cm}{0.4pt}
\end{center}
